\chapter{多项式：方程背后的结构}

\section{滑板公园的三个测量点}

城南新修的滑板公园里，有一座供滑板手腾空的抛台。设计师小周在图纸上画出了飞行轨迹：滑板手从台沿飞出，在空中划出一条弧线，再落回下坡的斜面上。这条弧线不是随手画的——它必须保证落点刚好在斜面最平缓的一段，否则轻则摔跤，重则出事故。

把场地搬进坐标系（这正是第 4 章学到的本领）：以起跳点为原点，水平方向为 $x$ 轴，竖直方向为 $y$ 轴，单位是米。小周知道，忽略空气阻力时，飞行轨迹的形状总是一条“二次曲线”，也就是可以用
\[
y = ax^2 + bx + c
\]
来描述的曲线，其中 $a$、$b$、$c$ 是三个待定的数。这类式子在生活中到处都是：喷泉水柱的弧线、篮球出手的路线、吊桥主缆的形状，背后都是同一个数学家族。麻烦在于：仪器只测到了轨迹上的三个点——起跳位置 $(0, 2)$，空中的一个点 $(1, 5)$，以及斜面上预先设定的安全落点 $(4, 2)$。

于是问题变成了两个。第一：有没有一条二次曲线同时穿过这三个点？如果有，它是唯一的吗？第二：假如轨迹确定下来，滑板手在什么水平位置时高度恰好是 $3$ 米？这相当于解方程 $ax^2+bx+c=3$，也就是问曲线和一条水平线在哪里相交。

先凭直觉猜一猜。你可能觉得三个条件定三个未知数，应该有解，而且多半只有一个——这个直觉是对的，但“为什么对”需要本章的工具才能讲清楚。你也可能觉得解方程 $ax^2+bx+c=3$ 无非是套求根公式——可如果算出来判别式是负的呢？那是“题目出错了”，还是数学在向我们暗示什么？

这一章的主线就从这两个疑问出发。我们把 $ax^2+bx+c$ 这样的式子当作一种有内部结构的“建筑”来研究，看看它的“零件”（因式）、“地基”（根）和“层数”（次数）如何决定它的全部行为；然后回到滑板公园，回答“三个点能不能还原一条规则”；最后直面那个负的判别式——它将把我们带到一个全新的数的世界。

\section{先试试直觉：两种容易掉的坑}

面对“曲线过三个点”，最自然的做法是凑。比如猜 $y = -x^2 + 4x + 2$：代入 $x=0$ 得 $y=2$，对了；代入 $x=1$ 得 $y=5$，又对了；代入 $x=4$ 得 $y = -16+16+2 = 2$，全对！运气不错，这回歪打正着。但如果换一组点，凑就难了——总不能每次都靠运气。我们需要一种“必能成功、且能证明唯一”的办法，本章后面会给出。

面对方程，直觉也容易踩坑。很多同学背熟了“图像与 $x$ 轴的交点就是方程的根”，于是顺推出两个看似合理的说法：

\begin{itemize}
\item 曲线碰了几次 $x$ 轴，方程就有几个根；
\item 图像不碰 $x$ 轴，方程就“没有意义”。
\end{itemize}

第一句会漏掉一种狡猾的情况：曲线只是“贴着”$x$ 轴反弹，并没有穿过去，比如 $y=(x-1)^2$ 在 $x=1$ 处轻轻碰一下就掉头。这时方程只有一个根，但这个根“占了两份”。第二句更危险：$x^2+1=0$ 的图像确实碰不到 $x$ 轴，可本章结尾会看到，它不但有意义，还逼出了数学史上最重要的一次“数的扩张”。

\begin{fanli}
“二次方程要么有两个根，要么没有实根。”——错。$x^2-2x+1=0$ 只有一个根 $x=1$，但它是“重复”的：左边就是 $(x-1)^2$，因子 $x-1$ 出现了两次。如果只说“一个根”，我们就丢掉了结构信息；更准确的说法是“一个二重根”。看图像时也请留心：与 $x$ 轴相切和穿过 $x$ 轴，是两种本质不同的相遇方式。
\end{fanli}

这两处失败告诉我们同一件事：光看“答案”不够，得拆开式子看结构。下面我们把“拆”这件事变成正式的工具。

\section{把式子当建筑：多项式、次数与根}

把 $3x^2-4x+1$ 拆开看，它由若干“块”相加而成：$3x^2$、$-4x$、$1$。每一块都是“一个数乘上 $x$ 的若干次方”（$1$ 可以看作 $1\cdot x^0$）。这样的式子用起来极其顺手：能相加、能相乘，乘完还是同一类式子，像乐高积木一样怎么搭都不出界。为了反复谈论这类对象，我们给它一个名字。

\begin{dingyi}[多项式]
形如
\[
a_n x^n + a_{n-1} x^{n-1} + \cdots + a_1 x + a_0
\]
的式子叫作（关于 $x$ 的）\textbf{多项式}，其中 $a_n, a_{n-1}, \ldots, a_0$ 是确定的数，叫作\textbf{系数}，$n$ 是非负整数。若 $a_n \neq 0$，则称这个多项式的\textbf{次数}为 $n$，$a_n$ 叫作\textbf{首项系数}。
\end{dingyi}

\begin{shuyu}{次数}
直观解释：多项式里“长得最快”的那一项的指数，决定了 $x$ 很大时整个式子的脾气。

正式定义：系数非零的最高次项的指数。

例子：$x^3-7x+6$ 的次数是 $3$；常数 $5$（即 $5x^0$）的次数是 $0$。
\end{shuyu}

\begin{shuyu}{根}
直观解释：让多项式“归零”的输入值；图像上就是曲线与 $x$ 轴相遇的位置。

正式定义：若数 $r$ 使 $f(r)=0$，则称 $r$ 是多项式 $f(x)$ 的一个根。

例子：$1$ 和 $3$ 都是 $x^2-4x+3$ 的根，因为代入后都得到 $0$。
\end{shuyu}

为什么要专门发明“次数”这个词？因为它像建筑物的层数：层数几乎决定了建筑的宏观特征。一次多项式的图像是直线，两端一头朝上一头朝下；二次的是抛物线；三次的像一条被扭了一下的丝带，从一端钻下去又从另一端冒出来。次数还回答了那个让古人纠结的问题——“方程最多有几个解”。

顺便说一句，多项式值得专门研究，还因为它“足够简单，又足够万能”。说它简单，是因为算它的值只需要乘法和加法，没有任何花哨操作，计算机最喜欢的正是这种式子；说它万能，是因为后面（第 11、12 章）我们会看到，连 $\sin x$、$\mathrm{e}^x$ 这些看上去毫不相干的函数，都能用“无限长的多项式”任意逼近。弄懂多项式，等于给整个函数世界打下地基。

而“根”把代数与几何焊在一起：解方程 $f(x)=0$，等价于找曲线 $y=f(x)$ 与 $x$ 轴的交点。同一个对象，一副代数面孔，一副几何面孔，我们可以随时换着看。

再看“拆”。$x^2-4x+3$ 可以写成 $(x-1)(x-3)$。拆开之后，根几乎自己跳出来：乘积为零，当且仅当某个因子为零，所以根是 $1$ 和 $3$。反过来，知道根是 $1$ 和 $3$，也立刻能写出因子 $x-1$ 与 $x-3$。这种“根与因子互相翻译”的本领，值得升格为一条定理。

\section{因式定理：根与因子的翻译机}

\begin{dingli}[因式定理]
设 $f(x)$ 是多项式，$a$ 是一个数。则 $a$ 是 $f(x)$ 的根，当且仅当 $x-a$ 是 $f(x)$ 的因式，即存在多项式 $q(x)$ 使得
\[
f(x) = (x-a)\,q(x).
\]
\end{dingli}

\begin{proof}
两个方向分别验证。

（如果 $x-a$ 是因式）假设 $f(x)=(x-a)q(x)$。把 $x=a$ 代入，右边第一个括号变成 $a-a=0$，零乘任何数都是零，所以 $f(a)=0$，即 $a$ 是根。

（如果 $a$ 是根）回忆小学学的带余除法：$17$ 除以 $5$ 商 $3$ 余 $2$，即 $17 = 5\times 3 + 2$，余数必须比除数小。多项式也有同样的除法：用 $f(x)$ 除以一次式 $x-a$，可以得到
\[
f(x) = (x-a)\,q(x) + r,
\]
其中“余数” $r$ 的次数必须低于除式 $x-a$ 的次数 $1$，所以 $r$ 只能是一个常数。现在把 $x=a$ 代入这个等式：左边是 $f(a)=0$（因为 $a$ 是根），右边第一项消失，剩下 $r$。于是 $r=0$，即 $f(x) = (x-a)q(x)$，$x-a$ 确实是因式。
\end{proof}

\begin{li}
验证 $x=2$ 是 $f(x)=x^3-7x+6$ 的根，并把它彻底拆开。代入：$f(2)=8-14+6=0$，是根。于是 $x-2$ 是因式。做除法（或耐心凑一凑）可得
\[
x^3-7x+6 = (x-2)(x^2+2x-3) = (x-2)(x-1)(x+3).
\]
三个根一目了然：$2$、$1$、$-3$。原来“解三次方程”这件事，一旦找到一个根，就降级成了“解二次方程”。
\end{li}

那一步除法具体怎么做？和整数竖式几乎一模一样。用 $x^3-7x+6$ 除以 $x-2$：先看最高次，$x^3$ 里藏着多少个 $x$ 倍的 $x-2$？商的第一项是 $x^2$，乘回去得 $x^3-2x^2$，从被除式里减掉，余下 $2x^2-7x+6$；再看 $2x^2$，商的下一项是 $2x$，乘回去得 $2x^2-4x$，减掉后余下 $-3x+6$；最后商 $-3$，乘回去得 $-3x+6$，刚好减干净，余数为 $0$。余数为零正是“$2$ 是根”的另一种写法。如果换一个不是根的数来除，最后总会剩下一小块——那一小块其实就是函数值，这就是继续攀登栏目里“余数定理”说的内容。

这个“降级”思想威力很大，而且可以反复使用，于是一个重要结论几乎免费到手：

\begin{tuilun}
$n$ 次多项式至多有 $n$ 个不同的实数根。
\end{tuilun}

\begin{proof}
每找到一个根 $a$，因式定理就送我们一个一次因式 $x-a$，多项式被拆成 $(x-a)q(x)$，其中 $q(x)$ 的次数降为 $n-1$。再找 $q(x)$ 的根，又能拆出一个一次因式。每拆一次，剩余部分的次数减一；最多拆 $n$ 次，次数就降到 $0$，剩下一个常数，再也没有根可拆。所以不同的根最多 $n$ 个。
\end{proof}

\begin{zhu}
“至多”二字不能省。$x^2+1$ 是二次的，却一个实根都没有。次数给的是根数的\emph{上限}，不是保证。这个缺口，本章结尾会回来补上。
\end{zhu}

\section{韦达关系：不拆也能听见根的声音}

拆因子很痛快，但有时拆不动（系数太丑），有时根本不需要拆到底。有没有办法不解方程，就知道根的一些信息？

对二次方程有，而且异常漂亮。设 $x^2+bx+c=0$ 的两个根是 $\alpha$ 和 $\beta$（允许相同）。由因式定理，
\[
x^2+bx+c = (x-\alpha)(x-\beta).
\]
把右边展开：
\[
(x-\alpha)(x-\beta) = x^2 - (\alpha+\beta)x + \alpha\beta.
\]
左右两边是同一个多项式，对应项的系数必须相同，于是我们得到：

\begin{dingli}[韦达关系（二次情形）]
若 $\alpha$、$\beta$ 是 $x^2+bx+c=0$ 的两个根，则
\[
\alpha+\beta = -b, \qquad \alpha\beta = c.
\]
对一般形式 $ax^2+bx+c=0$（$a\neq 0$），两边同除以 $a$ 化为上面的情形，可得 $\alpha+\beta=-\dfrac{b}{a}$，$\alpha\beta=\dfrac{c}{a}$。
\end{dingli}

\begin{li}
不解方程 $x^2-7x+12=0$：两根之和是 $7$，两根之积是 $12$。哪两个数和为 $7$、积为 $12$？$3$ 和 $4$。再看 $x^2+x+1=0$：如果它有实根，两根之和为 $-1$、积为 $1$。两个实数积为正，必须同号；和为负，必须同为负；可两个负数积为 $1$ 时，最接近的情况是 $-1$ 与 $-1$，和为 $-2$——比要求的 $-1$ 差了一点。这个“差一点”不是巧合，本章结尾会见分晓。
\end{li}

韦达关系真正的价值在于“隔空取物”：系数是根的\textbf{对称函数}。$\alpha+\beta$ 与 $\alpha\beta$ 有个共同特点——交换 $\alpha$ 和 $\beta$，它们的值不变。系数只能“看见”根的对称信息，看不见的部分（比如哪个根更大）它就说不出来。请记住“对称”这个词，它是通往下一章的钥匙之一。顺便一提，这个关系以十六世纪法国数学家韦达的名字命名；他的真正贡献不只是这两条公式，而是开始系统地用字母代表已知数和未知数——正是这套记号，让“研究所有方程”取代“解决一个个方程”成为可能。

三次、四次方程也有韦达关系：$x^3+px^2+qx+r=0$ 的三根之和为 $-p$，两两乘积之和为 $q$，三数之积为 $-r$。规律来自同样的展开，只是括号更多了。拿前面的老朋友 $x^3-7x+6$ 检验一下：它的三个根是 $1$、$2$、$-3$，和为 $0$（确实没有 $x^2$ 项），两两乘积之和为 $1\times 2+1\times(-3)+2\times(-3)=2-3-6=-7$（正是 $x$ 的系数），积为 $-6$（与常数项 $6$ 差一个负号，因为三次的常数项是“负的积”）。三个根的所有“对称信息”，被三个系数完整封存。

\section{图像是根的地图}

现在把代数得到的结构翻译成几何。面对一条多项式曲线，哪些信息是“读图可得”的？这一节给出的答案可以概括成一句话：图像不是别的，正是因式分解画出来的样子。我们分三点来看。

第一，\textbf{根在哪里，曲线就在哪里接触 $x$ 轴}，但接触的方式有讲究。$(x-1)(x-3)$ 在 $1$ 和 $3$ 处干脆地\emph{穿过}；$(x-1)^2$ 在 $1$ 处只是\emph{亲吻}一下就弹回。原因在于重数：因子 $x-a$ 出现奇数次时，$x$ 越过 $a$ 会让整个乘积变号，曲线穿轴；出现偶数次时乘积不变号，曲线触轴反弹。

\begin{center}
\begin{tikzpicture}[scale=0.9]
\draw[->] (-1,0) -- (5,0) node[right] {$x$};
\draw[->] (0,-1.8) -- (0,3.6) node[above] {$y$};
\draw[thick,domain=-0.3:4.3,smooth] plot (\x,{0.35*(\x-1)*(\x-3)});
\fill (1,0) circle (2pt) node[below] {$1$};
\fill (3,0) circle (2pt) node[below] {$3$};
\node at (2,1.1) {$y=(x-1)(x-3)$};
\end{tikzpicture}
\end{center}

\begin{center}
\begin{tikzpicture}[scale=1.0]
\draw[->] (-2.8,0) -- (2.8,0) node[right] {$x$};
\draw[->] (0,-2.4) -- (0,4.2) node[above] {$y$};
\draw[thick,smooth] plot coordinates {(-2.2,-2.0) (-2,0) (-1.5,3.1) (-1,4) (-0.5,3.4) (0,2) (0.5,0.6) (1,0) (1.5,0.9) (2,4)};
\fill (-2,0) circle (2pt) node[below left] {$-2$};
\fill (1,0) circle (2pt) node[below right] {$1$};
\node at (-1.9,-1.6) {穿过（单根）};
\node at (1.9,0.4) {反弹（二重根）};
\end{tikzpicture}
\end{center}

上面第二幅图画的是 $y=(x+2)(x-1)^2$：在 $-2$ 处穿过，在 $1$ 处反弹——图像把因式分解的秘密公开画了出来。

第二，\textbf{次数与首项系数决定两端的去向}。当 $x$ 的绝对值很大时，低次项全部变得微不足道，式子的行为由首项 $a_n x^n$ 一人说了算：$n$ 为偶数时两端同向（$a_n>0$ 都朝上，$a_n<0$ 都朝下）；$n$ 为奇数时两端反向。于是一个副产品自动出现：奇数次多项式一端冲上云霄、一端扎入地底，中间必然跨过 $x$ 轴——

\begin{tuilun}
奇数次多项式至少有一个实数根。
\end{tuilun}

严格证明需要“连续曲线不能跳过 $x$ 轴”这一事实（第 12 章的介值定理会把它变成铁律），但图像的直觉已经足够可信：从地底走到天空，不可能不穿过地面。

第三，\textbf{根把曲线切成若干段，每段内符号不变}。想知道 $x^3-7x+6>0$ 的解集？拆成 $(x+3)(x-1)(x-2)$，三个根 $-3$、$1$、$2$ 把数轴切成四段。每段取一个测试点判断正负即可，也可以列一张符号表：每个因子在自己的根右边为正、左边为负，整行的正负由负因子个数的奇偶决定。

\begin{center}
\begin{tabular}{ccccc}
\toprule
区间 & $x+3$ & $x-1$ & $x-2$ & 乘积 $f(x)$ \\
\midrule
$x<-3$ & $-$ & $-$ & $-$ & $-$ \\
$-3<x<1$ & $+$ & $-$ & $-$ & $+$ \\
$1<x<2$ & $+$ & $+$ & $-$ & $-$ \\
$x>2$ & $+$ & $+$ & $+$ & $+$ \\
\bottomrule
\end{tabular}
\end{center}

于是 $f(x)>0$ 的解集是 $-3<x<1$ 或 $x>2$，一行一行读表就行，不用再解任何方程。看上去麻烦的不等式，被“根地图”一次性解决。

\section{三个点猜一条规则：插值}

回到滑板公园：已知轨迹经过 $(0,2)$、$(1,5)$、$(4,2)$，求这条二次曲线。前面我们靠运气猜中了 $y=-x^2+4x+2$，现在给出必成之法，顺便把唯一性也证明了。

想法出人意料地简单：\textbf{先造“只在一点值班”的零件}。构造一个二次多项式 $\ell_1(x)$，要求它在 $x=0$ 处等于 $1$，在 $x=1$、$x=4$ 处等于 $0$。后两个条件用因式定理立刻办到：$\ell_1$ 必须含因子 $(x-1)(x-4)$，即 $\ell_1(x)=k(x-1)(x-4)$；再用“在 $0$ 处等于 $1$”定出常数 $k$：代入得 $k(0-1)(0-4)=4k=1$，所以 $k=\frac{1}{4}$，即
\[
\ell_1(x)=\tfrac{1}{4}(x-1)(x-4).
\]
同理，造出在 $x=1$ 处值班、其余两处为 $0$ 的零件 $\ell_2(x)=-\tfrac{1}{3}x(x-4)$（请验证：$\ell_2(1)=-\frac{1}{3}\cdot 1\cdot(-3)=1$），以及在 $x=4$ 处值班的零件 $\ell_3(x)=\tfrac{1}{8}x(x-1)$。

现在像调音台一样，把三个零件按目标高度混合：
\[
f(x) = 2\,\ell_1(x) + 5\,\ell_2(x) + 2\,\ell_3(x).
\]
检验：$x=0$ 时，$\ell_2$、$\ell_3$ 都为 $0$，$\ell_1=1$，所以 $f(0)=2$；同理 $f(1)=5$，$f(4)=2$。展开合并，$f(x)=-x^2+4x+2$——和猜的一样，但这次是算出来的，而且这套工序对任意三个点都有效。它有个名字：\textbf{拉格朗日插值}。

\begin{zhu}[唯一性]
过这三个点的二次多项式只有一个。假如 $f$、$g$ 都满足要求，考察它们的差 $f-g$：它在 $0$、$1$、$4$ 三处都为零——一个次数不超过 $2$ 的多项式竟然有三个根，这就违背了“二次至多两个根”的推论，除非 $f-g$ 恒等于零。看，刚证的根数上界立刻派上了用场。
\end{zhu}

一般规律是：$n+1$ 个横坐标不同的点，恰好确定一个次数不超过 $n$ 的多项式。点给少了，答案不唯一；点给多了，未必存在一个低次多项式刚好全部穿过。

这里藏着一个值得警惕的分寸：点多了，是不是次数越高越好？恰恰相反。假如测量有误差（现实里永远有），一条“硬要穿过每一个点”的高次曲线会在点与点之间疯狂扭动，预测能力反而糟糕。工程师们通常宁愿用一条低次曲线“大致穿过”所有点——这叫拟合，和插值是亲兄弟。插值负责“刚好经过”，拟合负责“总体最像”，分清什么时候用谁，是数据时代的基本功。

\begin{lianjie}
插值绝不只是几何游戏。手机放大照片时，新像素的颜色是插值算出来的；动画师只画关键帧，中间帧由软件插值生成；气象站只在若干城市有观测，全国气温图靠插值补全。更妙的是纠错码：它把数据看作某个多项式的函数值来传输，哪怕途中丢掉几个点，也能用插值把整条多项式还原出来——你听音乐、扫二维码时天天在用它。
\end{lianjie}

\section{当解不够用时：$i$ 的诞生}

现在来面对本章最大的那个疙瘩：$x^2+1=0$。

在实数范围内它无解——任何实数的平方都不是负数。几百年前的数学家也这么认为，并把这类方程扔进“无意义”的废纸篓。真正的转折发生在十六世纪的意大利：数学家们竞赛解三次方程，发现求根公式在中途要求对负数开平方。奇怪的是，如果\emph{假装}负数的平方根存在，硬着头皮按普通规则算下去，最后那些“虚”的部分竟然两两抵消，蹦出来的答案代回原方程——完全正确！

这就尴尬了：一个“不存在”的数，却在通往正确答案的路上充当了不可或缺的桥梁。拒绝它，明明存在的实根算不出来；接受它，一切顺畅。数学家们争论、怀疑了上百年——有人称这种数为“诡辩的量”，有人干脆拒绝书写它们——但计算的威力实在太诱人：凡是借道“虚数”算出的结果，最后都能用实数检验，而且次次正确。实践一再投票之后，理论终于让步。数学家们决定：把这座桥扶正，让它成为正式公民。规定一个新数 $i$，满足
\[
i^2 = -1.
\]
$i$ 不是实数轴上的居民——数轴上找不到它的位置。那它住在哪里？答案漂亮得让人释怀：既然实数占满了一条直线，就把 $i$ 放到直线\emph{外面}，放在平面上点 $(0,1)$ 处。于是一般的复数 $a+bi$ 就是平面上的点 $(a,b)$：实数躺在横轴上，$i$ 站在纵轴的单位处。

\begin{shuyu}{复数}
直观解释：平面上的一个点，横坐标是“实的部分”，纵坐标是“虚的部分”。

正式定义：形如 $a+bi$ 的数，其中 $a$、$b$ 是实数，$i$ 满足 $i^2=-1$。

例子：$3-2i$ 对应点 $(3,-2)$；实数 $5$ 就是 $5+0i$，纯虚数 $2i$ 就是 $0+2i$。
\end{shuyu}

复数的运算规则全部来自一条原则：像普通代数一样算，遇到 $i^2$ 就换成 $-1$。加法是对应分量相加：$(1+2i)+(3-i)=4+i$，在平面图上就是两个箭头首尾相接。乘法用分配律展开：
\[
(1+2i)(3-i)=3-i+6i-2i^2=3+5i-2\times(-1)=5+5i.
\]
每个复数 $a+bi$ 还有一位“镜像伙伴”$a-bi$，叫作它的\textbf{共轭}：两者相加得 $2a$（实数），相乘得 $a^2+b^2$（非负实数，恰好是到原点距离的平方）。靠这对伙伴，除法也能化为乘法——分子分母同乘分母的共轭，分母立刻变成实数。

\begin{center}
\begin{tikzpicture}[scale=1.4]
\draw[->] (-2.2,0) -- (2.4,0) node[right] {实轴};
\draw[->] (0,-1.6) -- (0,1.9) node[above] {虚轴};
\fill (1,0) circle (1.5pt) node[below right] {$1$};
\fill (-1,0) circle (1.5pt) node[below left] {$-1$};
\fill (0,1) circle (1.5pt) node[above right] {$i$};
\fill (0,-1) circle (1.5pt) node[below right] {$-i$};
\draw[->,thick] (0.75,0) arc (0:90:0.75);
\node at (0.95,0.7) {$\times i$};
\end{tikzpicture}
\end{center}

复数乘法有个几何解读：乘以 $i$ 相当于绕原点逆时针旋转 $90^\circ$。$1\times i = i$（从右方转到上方），$i\times i = i^2 = -1$（再转 $90^\circ$ 到左方）——所谓“负数的平方根”，原来就是“转半圈这个动作的平方根”，多么合理！

有了 $i$，$x^2+1=0$ 有了根 $\pm i$；$x^2+x+1=0$ 用求根公式得到
\[
x=\frac{-1\pm\sqrt{-3}}{2}=-\frac{1}{2}\pm\frac{\sqrt{3}}{2}i.
\]
还记得前面“差一点”的悬念吗？这两个复根的和不折不扣是 $-1$，积不折不扣是 $1$——韦达关系在复数世界里依然成立。更惊人的是下面这个事实：

\begin{dingli}[代数基本定理]
每个 $n$ 次（$n\ge 1$）复系数多项式在复数范围内至少有一个根；配合因式定理反复拆，它恰好可以拆成 $n$ 个一次因式的乘积，也就是说，恰有 $n$ 个根（重根按重数计数）。
\end{dingli}

这条定理的证明超出了本书的工具范围（它需要一点分析或拓扑的知识，大学的复变函数课会给出特别优雅的证法），但它的信息很清楚：\textbf{复数让“次数”和“根的个数”完美地对上了}。前面“奇数次必有实根”“二次可能没有实根”这类含糊的例外，在复数世界里全部消失——$n$ 次方程永远恰有 $n$ 个解，一个不多，一个不少。一种新数的引入，换来的不是更乱，而是惊人的整齐。

回头看，数的扩张走过一条熟悉的路：丈量土地不够用，有了分数；记账表示欠债不够用，有了负数；测量对角线不够用，有了无理数；解方程不够用，有了复数。每一次“解不够”，都不是数学的失败，而是数学在敲门：门后站着一个更大的数系，以及更整齐的定理。

\begin{pandeng}
继续向前，你会遇到这些词：\textbf{余数定理}（$f(a)$ 恰好等于 $f$ 除以 $x-a$ 的余数）、一般的\textbf{拉格朗日插值公式}、\textbf{共轭根定理}（实系数方程的虚根总是成对出现）、复平面上的\textbf{单位根}。再远处还站着一个大问题：“为什么五次及五次以上的方程没有通用求根公式？”它的答案诞生了一门用“对称”研究方程的学问——\textbf{群论}。
\end{pandeng}

\section{动手时间}

\begin{shiyan}
\textbf{实验一：根地图预测图像。}不看任何计算器，仅根据 $f(x)=(x+2)(x-1)(x-3)$ 回答：（1）曲线在哪三个位置接触 $x$ 轴？是穿过还是反弹？（2）$x$ 很大和很小时，两端各朝哪走？（3）在数轴上标出 $f(x)>0$ 和 $f(x)<0$ 的区间。然后描点画出草图，逐条核对你的预测。

\textbf{实验二：亲手做一次插值。}取三个点 $(0,1)$、$(2,3)$、$(3,1)$，仿照正文造出三个“值班零件”$\ell_1$、$\ell_2$、$\ell_3$，按高度混合出 $f(x)$，展开后把三个点代回检验。再任取一个横坐标（比如 $x=1$）算出函数值，描在坐标纸上，看看曲线是否真的光滑地串起了你的点。

\textbf{实验三：韦达侦探。}让一位同学在心里想好一个二次方程的两个整数根，只告诉你它们的和与积；你用韦达关系反推出方程，再求根验证。交换角色玩三轮，体会“系数只看得见对称信息”这句话的含义。
\end{shiyan}

\begin{toolbox}
本章新增的工具：
\begin{itemize}
\item \textbf{多项式、系数、次数}：把一大类式子当作统一的研究对象；
\item \textbf{因式定理}：根与一次因式互相翻译，找到一个根就把方程降一次；
\item \textbf{根数上界}：$n$ 次多项式至多 $n$ 个实根；奇数次至少有一个实根；
\item \textbf{韦达关系}：系数是根的对称函数，不解方程也能听到根的消息；
\item \textbf{重数与图像}：单根穿轴而过，重根触轴反弹；首项决定两端走向；
\item \textbf{拉格朗日插值}：$n+1$ 个点唯一确定一条不超过 $n$ 次的规则；
\item \textbf{复数}：$i^2=-1$；在复数世界里，$n$ 次方程恰有 $n$ 个解。
\end{itemize}
\end{toolbox}

\section*{思考题}

\textbf{观察题}
\begin{sikao}
\item 观察 $y=x^2$、$y=x^3$、$y=x^4$、$y=x^5$ 的草图。偶数次与奇数次多项式在 $x$ 很大和很小时各有什么规律？提示：只看首项，把其他项想象成“可以忽略的小波动”。
\item 画一画 $y=(x-1)^3$ 的草图。$x=1$ 是三重根，曲线在那里是穿过还是反弹？提示：奇数次方会改变符号，但增长在附近被“压平”了，描 $x=0.5$、$1$、$1.5$ 三点细看。
\end{sikao}

\textbf{证明题}
\begin{sikao}
\item 证明：若 $x^2+bx+c=0$ 有两个实根，则当两根之积为正时两根同号，为负时两根异号。提示：直接用韦达关系 $\alpha\beta=c$。
\item 已知 $f(x)$ 是三次多项式，且 $f(1)=f(2)=f(3)=0$，$f(0)=6$。求 $f(4)$。提示：先用因式定理写出 $f$ 的形状，只差一个待定系数。
\end{sikao}

\textbf{挑战题}
\begin{sikao}
\item 利用 $i^2=-1$ 计算 $(1+i)^2$ 和 $(1+i)^4$。提示：算完后把每个结果标在复平面上，观察每次平方时“到原点的距离”和“转角”发生了什么变化——这与“乘以 $i$ 等于转 $90^\circ$”暗暗呼应。
\item 证明：过 $n+1$ 个横坐标不同的点，至多只有一个次数不超过 $n$ 的多项式。提示：仿照正文中唯一性的注，考察两个候选多项式的差，数一数它有多少个根。
\end{sikao}

\section*{通往下一章的门}

这一章里，“对称”这个词悄悄出现了好几次：韦达关系里，$\alpha+\beta$ 与 $\alpha\beta$ 交换两个根而不变；系数永远只能“看见”根的对称信息；而那个终极问题——为什么三次、四次方程有求根公式，五次及以上却没有——最后的答案竟然是：\textbf{五次方程的根之间的“对称性”太复杂了，复杂到无法用有限次开方表达出来}。

给方程的根做“对称性体检”，听上去不可思议。但“对称”本身就是可以精确定义、可以运算的数学对象：正方形有几种翻折和旋转？这些动作怎样互相复合？先翻折再旋转，和先旋转再翻折，结果一样吗？下一章，我们把“对称”从形容词变成名词，从感觉变成结构——它会一路通向现代代数学的心脏：群。
